% Fresh strict and imported foundation checks passed; ordinary full checks are pending.
\section{A finite heat expansion of soft truncation}
\label{sec:soft_heat_expansion}

The finite logarithmic series is a finite linear combination of heat operators.
The coefficient budget and a logarithmic truncation order provide the scalar
estimates used in the corrected one-interval operator construction associated
with~\cite[lines 1262--1268]{Arad2016RigorousRG}.

\begin{definition}[Finite heat coefficients]
    \label{def:soft_truncation_heat_coefficients}
    \lean{Real.softTruncationHeatCoefficient}
    \leanok

    Define, for $0\leq a\leq m$,
    \begin{align}
        c_a(t,m)=t(-1)^a\sum_{j=1}^{m}\frac{\binom ja}{j}.
        \label{eq:soft_heat_coefficients}
    \end{align}
    Binomial coefficients with $a>j$ are zero. The constant
    coefficient is retained.
\end{definition}

\begin{theorem}[The exact scalar and matrix heat expansion]
    \label{thm:soft_truncation_heat_expansion}
    \lean{Real.softTruncation_eq_sum_heat,
        Matrix.IsHermitian.cfc_softTruncation_eq_sum_heat}
    \leanok

    \uses{def:soft_truncation_heat_coefficients}
    For $t>0$ and every real $x$,
    \begin{align}
        h_{m,t}(x)=\sum_{a=0}^m c_a(t,m)e^{-ax/t}.
        \label{eq:soft_heat_scalar}
    \end{align}
    For every finite Hermitian matrix $U$, the same identity gives
    \begin{align}
        h_{m,t}(U)=\sum_{a=0}^m c_a(t,m)e^{-aU/t}.
        \label{eq:soft_heat_matrix}
    \end{align}
\end{theorem}
\begin{proof}\leanok
    \uses{def:soft_truncation_heat_coefficients}
    Expand each $(1-e^{-x/t})^j$ by the binomial theorem, include
    the zero binomial terms to give all inner sums the same upper
    limit, and interchange the two finite sums. The identity
    $(e^{-x/t})^a=e^{-ax/t}$ identifies the heat times.
    Continuous functional calculus preserves these finite sums
    and scalar multiples.
\end{proof}

\begin{lemma}[Absolute coefficient budget]
    \label{lem:soft_truncation_heat_coefficient_budget}
    \lean{Real.sum_abs_softTruncationHeatCoefficient_le}
    \leanok

    \uses{def:soft_truncation_heat_coefficients}
    For $t\geq0$,
    \begin{align}
        \sum_{a=0}^m|c_a(t,m)|\leq t2^{m+1}.
        \label{eq:soft_heat_budget}
    \end{align}
\end{lemma}
\begin{proof}\leanok
    The unsigned binomial sum at degree $j$ is $2^j$.
    Hence the left side equals $t\sum_{j=1}^m2^j/j$, which is at
    most $t\sum_{j=1}^m2^j\leq t2^{m+1}$. This retains the
    factor $t$ in the coefficient budget.
\end{proof}

\begin{lemma}[A logarithmic order for an exponential geometric tail]
    \label{lem:exponential_geometric_tail_order}
    \lean{Real.exponential_geometric_tail_order}
    \leanok

    Let $W,\xi>0$, $0<b<1$, and $N\in\mathbb R$. Put
    \begin{align}
        \ell=\max\left\{1,
           \left\lceil\frac{\log\max\{1,N^2\}
                -\log\log(1+\xi/W)}{-\log b}\right\rceil\right\}.
        \label{eq:soft_heat_order}
    \end{align}
    Then $\ell\geq1$ and $W(e^{N^2b^\ell}-1)\leq\xi$.
\end{lemma}
\begin{proof}\leanok
    The ceiling gives
    $\log\max\{1,N^2\}+\ell\log b\leq\log\log(1+\xi/W)$.
    Exponentiating gives $N^2b^\ell\leq\log(1+\xi/W)$.
    Exponentiate again and multiply by $W$.
\end{proof}

\begin{theorem}[A polynomial bond bound at fixed filter degree]
    \label{thm:geometric_order_polynomial_bound}
    \lean{Real.exists_geometric_order_polynomial_bound}
    \leanok
    Fix $0<b<1$, $W,\xi>0$, and $q,s,k\in\mathbb N$ with $q\ge1$.
    There exist $C,p\in\mathbb N$, with $C\ge1$, such that every
    $N\in\mathbb N$ satisfies
    \begin{align}
        (2s r_N^2 q^{r_N})^{2k}&\le C(N+1)^p,\\
        r_N&=\max\left\{1,
          \left\lceil\frac{\log\max\{1,N^2\}
             -\log\log(1+\xi/W)}{-\log b}\right\rceil\right\}.
    \end{align}
    The coefficient and degree are chosen before $N$; zero length
    and $s=0$ or $k=0$ are included.
\end{theorem}
\begin{proof}\leanok
    Put $d=-\log b>0$, $L=\log\log(1+\xi/W)$,
    $a_0=2+|L|/d$ and $a_1=2/d$. With $X=N+1\ge1$, the ceiling
    estimate gives
    \begin{align}
        r_N\le a_0+a_1\log X\le(a_0+a_1)X.
    \end{align}
    For a negative argument the natural ceiling is zero; for a
    nonnegative argument it is at most the argument plus one.
    Also $\log\max\{1,N^2\}\le2\log X$.
    Put $p_0=\lceil a_1\log q\rceil$. Since $q\ge1$,
    \begin{align}
        q^{r_N}
          &\le e^{a_0\log q}X^{a_1\log q}
           \le e^{a_0\log q}X^{p_0},\\
        2s r_N^2q^{r_N}
          &\le KX^{p_0+2},\qquad
          K=2s(a_0+a_1)^2e^{a_0\log q}.
    \end{align}
    Take $C=\max\{1,\lceil K^{2k}\rceil\}$ and $p=(p_0+2)2k$.
    Raising the last inequality to the natural power $2k$ proves the claim.
\end{proof}

This numerical consequence applies to a fixed filter degree and fixed
approximation budgets. It gives no assertion for a filter degree that
grows with the length and no running-time estimate.
